\documentclass[12pt, a4paper]{article}

\usepackage[utf8]{inputenc}
\usepackage{graphicx}
\usepackage{hyperref}
\usepackage{geometry}
\usepackage{listings}
\usepackage{color}

\geometry{a4paper, margin=1in}

% Code snippet styling
\definecolor{codegreen}{rgb}{0,0.6,0}
\definecolor{codegray}{rgb}{0.5,0.5,0.5}
\definecolor{codepurple}{rgb}{0.58,0,0.82}
\definecolor{backcolour}{rgb}{0.95,0.95,0.92}

\lstdefinestyle{mystyle}{
    backgroundcolor=\color{backcolour},   
    commentstyle=\color{codegreen},
    keywordstyle=\color{magenta},
    numberstyle=\tiny\color{codegray},
    stringstyle=\color{codepurple},
    basicstyle=\ttfamily\footnotesize,
    breakatwhitespace=false,         
    breaklines=true,                 
    captionpos=b,                    
    keepspaces=true,                 
    numbers=left,                    
    numbersep=5pt,                  
    showspaces=false,                
    showstringspaces=false,
    showtabs=false,                  
    tabsize=2
}
\lstset{style=mystyle}

\title{\textbf{Document Verification System Using Blockchain}}
\author{College Project Report}
\date{\today}

\begin{document}

\maketitle
\tableofcontents
\newpage

\section{Abstract}
In the digital age, document forgery and unauthorized modifications are significant security threats. Traditional centralized verification systems are prone to single points of failure, data breaches, and tampering. This project presents a decentralized Document Verification System leveraging Blockchain technology. By storing cryptographic hashes of documents on an immutable ledger rather than the documents themselves, the system ensures data privacy while providing an indisputable proof of existence and authenticity. The project is implemented using Ethereum smart contracts (Solidity), local blockchain environment (Ganache), and a modern Web3 frontend interacting via Ethers.js.

\section{Introduction}
Document verification is a critical process in various domains, including education, legal, real estate, and finance. Current verification processes often require manual intervention or reliance on centralized authorities, which can be slow, costly, and vulnerable to corruption. 

Blockchain technology inherently provides immutability, transparency, and decentralization. By utilizing a blockchain network, we can create a trustless environment where anyone can verify the authenticity of a document without relying on a central server. This project implements a simple yet highly effective proof-of-concept for such a system.

\section{Objectives}
The primary objectives of this project are:
\begin{itemize}
    \item To develop a decentralized application (DApp) that allows users to register documents securely.
    \item To ensure document privacy by hashing files locally in the browser and only storing the hash on the blockchain.
    \item To provide a public verification mechanism where any user can upload a document to verify its authenticity and registration timestamp.
    \item To build a modern, intuitive, and user-friendly web interface for seamless interaction with the blockchain.
\end{itemize}

\section{System Architecture}
The system architecture follows a standard Web3 model, comprising the following components:

\subsection{Frontend Application}
The user interface is built using HTML, CSS (with modern glassmorphism design), and Vanilla JavaScript, powered by the Vite build tool. The frontend handles:
\begin{itemize}
    \item Connecting to the user's Web3 wallet (MetaMask).
    \item Using the Web Crypto API to generate a SHA-256 hash of the uploaded document entirely on the client side.
    \item Interacting with the smart contract via the Ethers.js library.
\end{itemize}

\subsection{Smart Contract}
The core logic resides in a Solidity smart contract named \texttt{DocumentVerifier.sol}. It maintains a mapping of document hashes to a struct containing the owner's address, a registration timestamp, and a boolean flag indicating existence. 

\subsection{Blockchain Network (Ganache)}
For development and demonstration purposes, the smart contract is deployed on Ganache, a personal local Ethereum blockchain. Hardhat is utilized as the development environment for compiling, testing, and deploying the contract to Ganache.

\section{Implementation Details}

\subsection{Cryptographic Hashing for Privacy}
A crucial design decision in this system is that \textbf{actual documents are never uploaded to any server or blockchain}. When a user selects a file, the frontend JavaScript reads the file locally and computes its SHA-256 hash. Only this 64-character hexadecimal string is sent in the blockchain transaction. This guarantees absolute privacy for sensitive documents.

\subsection{Smart Contract Logic}
The smart contract contains two primary functions:
\begin{itemize}
    \item \textbf{registerDocument(string memory \_hash)}: Takes the document hash, checks if it already exists to prevent duplicates, and maps it to the sender's address (\texttt{msg.sender}) and the current block timestamp.
    \item \textbf{verifyDocument(string memory \_hash)}: A read-only function that returns the owner and timestamp of a given hash if it exists in the registry.
\end{itemize}

\subsection{User Interface \& Experience}
The application features a dark-themed, glassmorphism UI for a premium look. It provides drag-and-drop zones for both registering and verifying documents. The UI handles errors gracefully, such as duplicate registrations or attempting to interact without a connected wallet, providing immediate visual feedback to the user.

\section{Testing and Deployment}
The system was tested using the Hardhat framework. The deployment process involves:
\begin{enumerate}
    \item Starting the local Ganache RPC server on port 7545.
    \item Compiling the Solidity contract using \texttt{npx hardhat compile}.
    \item Deploying the contract via \texttt{npx hardhat run scripts/deploy.js -{}-network ganache}.
    \item Automatically exporting the generated ABI and Contract Address to the frontend source files.
\end{enumerate}

\section{Conclusion}
This project successfully demonstrates the viability of blockchain technology for secure, decentralized document verification. By combining client-side cryptographic hashing with Ethereum smart contracts, the system eliminates the need for trusted third parties while ensuring the privacy and immutability of document records.

\section{Future Scope}
While the current implementation is a robust proof-of-concept, future enhancements could include:
\begin{itemize}
    \item Deploying the smart contract to a public testnet (e.g., Sepolia) or a Layer 2 scaling solution (e.g., Polygon) for real-world usability and lower gas fees.
    \item Adding a revocation feature, allowing document owners to invalidate a previously registered document.
    \item Integrating IPFS (InterPlanetary File System) for optional decentralized storage of the actual documents for public records.
    \item Multi-signature approval for enterprise document registration workflows.
\end{itemize}

\end{document}
